% !TeX root = ../教材题库_全部条目_转写工作稿.tex
% Source: XJ-B1-C06-HW-P231-Y0872
% 第6章 / 6.3 统计图表 / 教材 p.231 / PDF 第238页
% OCR未校对：公式与文字均可能有误；未完成数学TeX转写；禁止直接发给学生
% 原题及相邻原解截图（核对用）：../题目原图/XJ-B1-C06-HW-P231-Y0872-01.webp
% 本文件只在转写完成后作为题干使用；OCR 条目尚不是可靠题干。
\begin{bankitem}{习题6.3 第2题}{第6章 / 6.3 统计图表 / 教材 p.231 / PDF 第238页}{XJ-B1-C06-HW-P231-Y0872}{OCR未校对：公式与文字均可能有误；未完成数学TeX转写；禁止直接发给学生}
\begin{ocrdraft}
[教材p.231 / PDF第238页]
\par
2．20世纪初，人们将驼鹿引入美国密歇根湖的一座孤岛．
\par
该种群从1915年到
\par
1960年的数量变化情况如下表：
\par
19巧19171921192519281930
\par
1934
\par
1943
\par
1947
\par
1950
\par
1960
\par
驼鹿种群数量/只2003001000200025003000
\par
00
\par
170
\par
600
\par
500
\par
600
\par
(1)用统计图表示该种群数量随时间变化的情况．
\par
（2）从1915年到1930年，该种群数量不断增加，可能的原因有哪些？
\par
（3）该种群的数量后来急剧下降，可能的原因有哪些？
\end{ocrdraft}
\end{bankitem}
